BitBucket Cloud uses the top-level concept of a \textbf{Workspace} to logically group multiple
related repositories.  A BitBucket tenant with a premium subscription may have multiple workspaces
but a non-premium tenant will have at most one workspace.

A workspace contains one or more \textbf{Project} configuration.  Creating a \textbf{Repository}
requires that it is assigned to a project.  The names of created repositories must be unique
in the workspace regardless of project assignment.

Webhooks can be deployed at either the \textbf{Workspace} or individual \textbf{Repository} scope.
Deploying webhooks at the \textbf{Workspace} scope sends events to \cxoneflow for all repositories
in the workspace.  While there is a guided UI for deploying repository webhooks, deploying webhooks
at the workspace scope is available only by using the BitBucket Cloud REST API.  The \cxoneflowaudit
tool can be used to deploy the webhook configuration at the workspace scope.

Webhooks can be deployed at the scope of each \textbf{Repository} if desired.  The number of repositories
in a large enterprise generally makes deployment at the \textbf{Repository} scope useful
only for testing purposes.
